% Source: https://artofproblemsolving.com/community/c3h1916761_math_problems
% Stable identity: aops-topic-1916761-q03
\textbf{All regions.} When the $k$th line is added, it meets the previous $k-1$ lines in $k-1$ different points, which cut it into $k$ pieces; each piece splits one existing region into two. So $n$ lines make
\[
1+(1+2+\cdots+n)=1+n+\dbinom n2
\]
regions; for $n=17$ this is $1+17+136=154$.

\textbf{Unbounded regions.} Draw a circle so large that all intersection points lie inside it. Outside the circle the lines are $2n$ disjoint rays, which cut the outside into $2n$ pieces; every unbounded region contains exactly one of these pieces, and the bounded regions lie inside the circle. So there are $2n=34$ unbounded regions.

Hence the number of closed regions is $154-34=\boxed{120}$ $\left(=\dbinom{16}{2}\right)$.
